\documentclass[10pt]{article}
\usepackage[a4paper, margin=1.9cm]{geometry}
\usepackage{booktabs}
\usepackage{parskip}
\usepackage{titlesec}
\titlespacing*{\section}{0pt}{8pt}{3pt}
\titleformat{\section}{\normalfont\large\bfseries}{\thesection}{0.6em}{}
\pagenumbering{gobble}

\title{\vspace{-2cm}A Hypothesis-Driven CHIA Loop for Cross-SoC Cache Hierarchy Optimization}
\date{\vspace{-0.5cm}}

\begin{document}
\maketitle
\vspace{-1cm}

\section{Task Overview}
This proposal addresses the cache hierarchy optimization track. The cache hierarchy can be one of the highest leverage design choices in an SoC, yet tuning it means searching a space of sizes, associativities, line sizes, replacement policies, and prefetchers far too large to simulate exhaustively. Worse, this search restarts from zero for every new SoC, while a human architect carries understanding between designs. Existing agentic approaches largely treat optimization as black-box search: they often require large simulation budgets and typically do not preserve explicit, testable architectural reasoning that can be transferred to the next design. We propose a CHIA loop in which the LLM agent works the way an architect does: it reads simulation results, states a hypothesis about the current bottleneck, and runs only the experiments that test it, while a statistical surrogate model discards experiments whose outcome is already predictable. We run this loop on three heterogeneous SoCs to test our central hypothesis: explicit, hypothesis-level architectural knowledge transfers across SoCs more effectively than raw observations or a surrogate model alone, cutting the simulations required to optimize an unseen design.

\section{Methodology and Loop Design}
\textbf{Setup.} We use ChampSim, wrapped as a CHIA tool node and parallelized on GCP, and instantiate three deliberately heterogeneous SoC profiles (small mobile-like, mid-range, and server-class) with materially different cache capacities, latencies, and core and memory-system parameters, each evaluated on a fixed mix of SPEC CPU2017 traces. The knobs are L1/L2/LLC size, associativity, line size, replacement policy, and prefetcher; the objective is IPC under an area proxy constraint.

\textbf{The loop.} A small seed batch of simulations initializes the loop. Each round, an analyst agent (Gemini) reads the accumulated results as compact metric tables and writes a hypothesis about the current bottleneck with candidate configurations to test it. A surrogate model trained on all simulations so far scores candidates by predicted utility and uncertainty, suppressing redundant experiments while retaining uncertain or high-value ones; CHIA runs the selected candidates in parallel. The hypothesis is accepted or rejected against the new evidence and the round is logged to a machine-generated research notebook. The loop stops when the budget is spent or improvement stalls.

\textbf{Transfer phase.} After converging on SoC A and B, the agent distills its notebook into explicit rules and optimizes the unseen SoC C under a hard simulation cap. Unlike Bayesian optimization, which transfers only through a statistical model of the design space, these rules (e.g. whether a workload is capacity-, conflict-, or latency-limited) remain valid priors even when optimal parameter values differ. To isolate the source of transfer gains, we evaluate (i) no transfer, (ii) pooled-surrogate transfer, (iii) distilled-rule transfer, and (iv) their combination.

\textbf{Execution plan.} Week 1: ChampSim node, traces, reference sweeps for ground truth. Week 2: surrogate and agent nodes, single-SoC loop end to end. Week 3: runs on A and B, transfer and ablations on C. Week 4: analysis, open-source cleanup, paper.

\section{Expected Results}
(1) \emph{Sample efficiency:} we test whether the hybrid loop can reach within 2\% of the best configuration found by the dense reference sweep using substantially fewer simulations than random and surrogate-only search.
(2) \emph{Transfer:} on the unseen SoC C, we measure the reduction in simulations required to reach fixed IPC-quality thresholds, comparing rule transfer, surrogate transfer, their combination, and scratch baselines.
(3) \emph{Explanation:} a research notebook per SoC tying each accepted change to a tested hypothesis.
(4) \emph{Artifacts:} an open-source, reproducible CHIA loop whose ChampSim wrapper, surrogate, and analyst nodes are reusable as standalone CHIA blocks.

\section{Cost Estimate}
A ChampSim run takes roughly 10 minutes on one CPU core; we budget about 5{,}000 simulations across reference sweeps, agent loops, and baselines.

\begin{center}
\begin{tabular}{lr}
\toprule
GCP compute ($\sim$1{,}000 core-hrs, incl. headroom for reruns) & \$260 \\
Gemini API (analyst calls over 4 weeks) & \$290 \\
Storage, orchestration VM, misc. & \$50 \\
\midrule
\textbf{Total} & \textbf{$\sim$\$600} \\
\bottomrule
\end{tabular}
\end{center}

This sits within the \$100 to \$1{,}500 target, with room to scale up the trace mix or add a fourth SoC.

\end{document}
